\documentclass[11pt,a4paper]{article}
\usepackage[T1]{fontenc}
\usepackage[utf8]{inputenc}
\usepackage[main=italian,provide=*]{babel}
\usepackage{amsmath,amssymb}
\usepackage{geometry}
\geometry{margin=2.3cm,top=2.5cm,bottom=2.7cm}
\usepackage{booktabs}
\usepackage[table,dvipsnames]{xcolor}
\usepackage{tcolorbox}
\tcbuselibrary{skins,breakable}
\usepackage{titlesec}
\usepackage{enumitem}
\usepackage{needspace}
\usepackage{fancyhdr}
\usepackage{graphicx}
\usepackage{hyperref}
\hypersetup{colorlinks=true,linkcolor=Sepia,citecolor=Sepia,urlcolor=Sepia}
\newcommand{\Tr}{\operatorname{Tr}}
\newcommand{\ket}[1]{\lvert #1 \rangle}
\newcommand{\bra}[1]{\langle #1 \rvert}

\definecolor{inkblue}{HTML}{1B3A5C}
\definecolor{lightblue}{HTML}{EAF1F8}
\definecolor{accent}{HTML}{B0562A}
\definecolor{softgray}{HTML}{6B6B6B}
\definecolor{okgreen}{HTML}{2E6B3E}
\definecolor{okgreenbg}{HTML}{EAF5EC}

\titleformat{\section}
  {\normalfont\Large\bfseries\color{inkblue}}
  {\thesection}{0.6em}{}[{\color{inkblue!40}\titlerule}]
\titlespacing{\section}{0pt}{0.9em}{0.4em}
\titleformat{\subsection}
  {\normalfont\large\bfseries\color{accent}}
  {}{0em}{}
\titlespacing{\subsection}{0pt}{0.6em}{0.2em}

\pagestyle{fancy}
\fancyhf{}
\renewcommand{\headrulewidth}{0.4pt}
\fancyhead[L]{\small\color{softgray}Scan sui parametri di rumore (anello)}
\fancyhead[R]{\small\color{softgray}\thepage}
\fancyfoot[C]{\small\color{softgray}Tesi triennale, Samuele --- Relatore: Prof.\ Alessandro Chiesa}

\newtcolorbox{keyeq}[1][]{
  colback=lightblue, colframe=inkblue!70, boxrule=0.6pt,
  arc=2pt, left=8pt, right=8pt, top=3pt, bottom=3pt, #1
}
\newtcolorbox{okbox}[1]{
  colback=okgreenbg, colframe=okgreen!60, boxrule=0.6pt,
  arc=2pt, left=7pt, right=7pt, top=2pt, bottom=2pt,
  fonttitle=\bfseries\color{okgreen}, title={#1}
}

\title{\vspace{-1.5em}\color{inkblue}Scan sui parametri di rumore --- trimero ad anello\\[0.2em]
\Large\normalfont\color{softgray}Come si sposta $N^*$, in entrambi gli scenari}
\author{Note di lavoro --- Tesi triennale, Samuele\\ Relatore: Prof.\ Alessandro Chiesa}
\date{}

\begin{document}
\sloppy
\maketitle
\thispagestyle{fancy}
\vspace{-2.0em}

File di codice: \texttt{scan\_\allowbreak parametri\_rumore\_trimero\_anello.py}
(mirror di \texttt{scan\_parametri\_rumore\_dimero.py}, esteso a due
scenari). Metrica: fedeltà di Trotter (Stadio 3, osservabile più
economico), entrambi gli scenari A (VQE) e B ($R_0$). Invarianza da
$p_\text{readout}$ verificata sul correlatore (Stadio 4).

\tableofcontents

\section{Obiettivo}

Gli stadi precedenti hanno trovato $N^*$ a un \textbf{singolo} punto
della calibrazione di riferimento (\texttt{ibm\_torino}): $N^*_A=3$,
$N^*_B=9$ (fedeltà, Stadio 3, valori corretti --- Sez.~\ref{sec:correzione}),
$N^*=3$ per il correlatore (Stadio 4, stesso punto dello scenario A).
Questo documento chiede: $N^*$ si sposta se il rumore è più o meno
forte? E se il readout cambia? A differenza del dimero, qui la domanda
va posta per \textbf{due} scenari fisici distinti, non uno solo.

\begin{keyeq}[title={Previsione, prima della verifica --- stessa del dimero}]
Rumore di gate ($\varepsilon_{1q},\varepsilon_{2q}$): ci si aspetta
$N^*$ \textbf{decrescente}. Readout ($p_\text{readout}$): fattore
puramente moltiplicativo $(1-2p_\text{readout})$ su ogni $N$ allo stesso
modo (Stadio 4) --- non sposta mai un argmax, quindi $N^*$ atteso
\textbf{invariante}.
\end{keyeq}

\section{Correzione: i valori di \texorpdfstring{$N^*$}{N*} dello Stadio 3 erano da una griglia troppo grezza}
\label{sec:correzione}

\begin{okbox}{Correzione, propagata anche agli Stadi 3 e 4}
Nel costruire lo scan qui è emerso (Sez.~\ref{sec:regressione}) che la
griglia $N\in\{1,2,5,10,\ldots\}$ usata allo Stadio 3 saltava $N=3$ e
$N=4$: i veri massimi, con griglia continua $N=1,\ldots,20$, sono
$N^*_A=3$ (non $2$) e $N^*_B=9$ (non $10$). Corretto retroattivamente
nei documenti degli Stadi 3 e 4 (figura e testo rigenerati), non solo
qui. Conseguenza non banale: $N^*_A=3$ \textbf{coincide} con $N^*$ del
correlatore (Stadio 4, stesso punto di lavoro) --- prima sembravano
tre valori diversi, ora sono due.
\end{okbox}

\section{Risultato: nessuna regressione, griglia fine confermata}
\label{sec:regressione}

\begin{center}
\renewcommand{\arraystretch}{1.25}
\begin{tabular}{@{}lcc@{}}
\toprule
\textbf{Scenario} & $\boldsymbol{N^*}$ \textbf{(griglia fine)} & $\boldsymbol{F(N^*)}$\\
\midrule
A (VQE) & 3 & 0.700335\\
B ($R_0$) & 9 & 0.381896\\
\bottomrule
\end{tabular}
\end{center}

\begin{okbox}{Verificato}
Ricalcolato con \texttt{trova\_N\_star} (funzione generale, cammino di
codice diverso dalla ricerca manuale degli Stadi 3-4): stessi valori,
coincidenza esatta.
\end{okbox}

\section{Risultato: monotonia rispetto a \texorpdfstring{$\varepsilon_{1q},\varepsilon_{2q}$}{eps}, entrambi gli scenari}

\begin{figure}[h]
\centering
\includegraphics[width=0.98\linewidth]{figure/fig_Nstar_vs_eps_trimero_anello.pdf}
\caption{\textit{$N^*$ in funzione di $\varepsilon_{2q}$ (sinistra) ed
$\varepsilon_{1q}$ (destra), scala log-log, l'altro parametro fissato al
riferimento. Linea tratteggiata grigia: valore di calibrazione
\texttt{ibm\_torino}. Monotonia non-crescente confermata in entrambi gli
scenari e in entrambi i parametri, con saturazione a $N^*=1$ per rumore
sufficientemente forte --- lo scenario B satura più bruscamente
(transizione netta fra $\varepsilon_{2q}=6\times10^{-3}$ e
$10^{-2}$), lo scenario A più gradualmente.}}
\end{figure}

\begin{okbox}{Verificato}
$N^*$ monotono non-crescente in $\varepsilon_{1q}$ e $\varepsilon_{2q}$,
per \textbf{entrambi} gli scenari --- non assunto che valesse per uno
solo. La saturazione a $N^*=1$ conferma lo stesso meccanismo del dimero:
oltre una certa soglia di rumore, nessun $N>1$ vale il costo aggiuntivo
in gate rumorosi.
\end{okbox}

\section{Risultato: invarianza rispetto a \texorpdfstring{$p_\text{readout}$}{p readout}}

\begin{center}
\renewcommand{\arraystretch}{1.2}
\begin{tabular}{@{}ccc@{}}
\toprule
$\boldsymbol{p_\text{readout}}$ & $\boldsymbol{N^*}$ (correlatore) & $|\boldsymbol{C(N^*)}|$\\
\midrule
$0.00$ & $3$ & $0.5292$\\
$0.05$ & $3$ & $0.4763$\\
$0.20$ & $3$ & $0.3175$\\
\bottomrule
\end{tabular}
\end{center}

\begin{okbox}{Verificato: la previsione teorica regge}
$N^*=3$ per tutti i valori di $p_\text{readout}$ testati, dallo $0\%$ al
$20\%$. Il valore assoluto $|C(N^*)|$ cambia sensibilmente, la
\emph{posizione} del massimo no --- firma di un fattore moltiplicativo
costante in $N$, stessa proprietà verificata sul dimero e già
argomentata (non solo assunta) allo Stadio 1.
\end{okbox}

\section{Sul modello additivo: nessuna ragione di aspettarsi che funzioni qui}

Sul dimero, il modello ingenuo $\varepsilon_\text{tot}(N)\sim
a/N+g\varepsilon_{2q}N$ prevedeva una pendenza log-log di $-0.5$ per
$N^*(\varepsilon_{2q})$; osservata invece $-0.195$ --- dichiarato
esplicitamente nel documento del dimero come mismatch atteso (la
fedeltà si comporta più come un decadimento moltiplicativo che come una
somma di due termini). Nessuna ragione di aspettarsi un accordo migliore
qui: due osservabili di fedeltà (scenari A e B) più un correlatore, tutte
governate dallo stesso meccanismo di rumore depolarizzante accumulato.
Non è stato quindi eseguito un fit di pendenza per l'anello --- sarebbe
un esercizio di conferma di un disaccordo già capito, non un'informazione
nuova.

\section{Cosa non è coperto da questo scan}

Un parametro alla volta, con gli altri fissati al riferimento --- non
l'intero spazio congiunto $(\varepsilon_{1q},\varepsilon_{2q},
p_\text{readout})$. Il correlatore è stato scansionato solo per
l'invarianza da readout (Sez.~4), non per $\varepsilon_{1q}/\varepsilon_{2q}$
--- per costo computazionale, stessa scelta motivata del dimero
(l'osservabile più economico, la fedeltà, basta a rispondere alla
domanda "come si sposta $N^*$" per il rumore di gate).

\section{Considerazioni finali sulla pipeline di rumore}

Questo chiude la pipeline in cinque stadi sul trimero ad anello: dal
modello di rumore (wrapper sul dimero, Stadio 1) fino a qui, ogni stadio
ha aggiunto un livello di struttura --- energia statica (Stadio 2), poi
dinamica in due scenari fisici distinti (Stadio 3), poi un osservabile a
due tempi con misura reale (Stadio 4), poi la sensibilità del risultato
ai parametri del rumore stesso (qui). Il filo conduttore verificato ad
ogni stadio: il readout simmetrico è sempre un fattore moltiplicativo,
mai additivo.

Una correzione di percorso, non nascosta: una griglia $N$ troppo
grezza allo Stadio 3 aveva riportato $N^*_A=2$ invece di $3$
(Sez.~\ref{sec:correzione}) --- scoperta solo costruendo lo scan di
questo stadio con una griglia più fitta, poi propagata a ritroso. Lo
stesso principio già applicato più volte in questa sessione: un
controllo diretto batte un'estensione per analogia, e un numero riportato
resta provvisorio finché non è ricontrollato con una risoluzione
sufficiente.

\begin{okbox}{Sintesi in una riga}
$N^*$ dipende in modo monotono e prevedibile dal rumore di gate, in
entrambi gli scenari fisici; non dipende affatto dal readout simmetrico
--- entrambi i fatti verificati, non solo previsti; e la griglia $N$
usata per trovarlo deve essere abbastanza fine da non saltare il vero
massimo, lezione imparata a proprie spese in questo stadio.
\end{okbox}

\needspace{4\baselineskip}
\section*{File prodotti}
\addcontentsline{toc}{section}{File prodotti}

\texttt{scan\_\allowbreak parametri\_rumore\_trimero\_anello.py},
\texttt{validate\_\allowbreak scan\_\allowbreak parametri\_rumore\_trimero\_anello.py},
\texttt{generate\_figure\_\allowbreak Nstar\_vs\_eps\_trimero\_anello.py},
questo documento. File corretti retroattivamente:
\texttt{risultati\_\allowbreak trotter\_\allowbreak rumoroso\_trimero\_anello.tex/.pdf},
\texttt{risultati\_\allowbreak correlatori\_\allowbreak rumorosi\_trimero\_anello.tex/.pdf}
(Sez.~\ref{sec:correzione}).

\end{document}
